% Fully read excerpts from https://raw.githubusercontent.com/openai/math/fd4aeeb2ee4fc729c18d98444fed42fd0529eeeb/preprints/A-direct-proof-of-the-complete-Crouzeix-inequality-September-26-2026/build/main.tex

\begin{theorem}\label{thm:main}
For every pair of positive integers $n,m$, every $A\in M_n(\C)$,
every nonnegative integer $d$, and every polynomial
$P(z)=\sum_{k=0}^d B_kz^k$ with $B_k\in M_m(\C)$,
\begin{equation}\label{eq:main}
 \|P[A]\|\le2\max_{z\in W(A)}\|P(z)\|.
\end{equation}
The constant two cannot be decreased uniformly over $n,m$, and $d$.
\end{theorem}

The estimate includes point and segment numerical ranges. Its complete
character is the uniformity in $m$: a scalar proof need not survive
matrix amplification, because coefficient matrices need not commute.
The argument below keeps their multiplication order, using the two
products $FG$ and $GF$ separately.



% --- next excerpt ---

\section{Exterior coordinates and coefficient estimates}\label{sec:coefficients}
In this section $\Omega$ is a bounded open convex set with regular real-analytic Jordan boundary. We work with a supplied exterior conformal map $h$ that is one-to-one on $|\lambda|>r$, including infinity on the Riemann sphere, for some $0<r<1$. It maps $|\lambda|>1$ onto $\C\setminus\overline\Omega$ and has a simple pole at infinity, with expansion
\begin{equation}\label{eq:h}
 h(\lambda)=c\lambda+c_0+O(\lambda^{-1}),\qquad c\ne0.
\end{equation}
Appendix~\ref{sec:collar} supplies such a coordinate for every domain under consideration. Here we use it to compare polynomial coefficients with boundary norms.

The map $h$ sends $\T=\{\lambda:|\lambda|=1\}$ onto $\partial\Omega$. Put $q(\lambda)=\lambda h'(\lambda)$. On $\T$, the vector $q(\lambda)$ is an outward normal and $iq(\lambda)$ is a positively oriented tangent. Hence
\begin{equation}\label{eq:support}
 \Re\bigl(\overline{q(\lambda)}(h(\lambda)-z)\bigr)\ge0
 \quad(\lambda\in\T,\ z\in\overline\Omega).
\end{equation}
Figure~\ref{fig:normal} illustrates this supporting-line relation.

\input{figures/normal.tex}

Every circle integral below uses $d\tau(e^{it})=dt/(2\pi)$.

We use the unnormalized Hilbert--Schmidt norm $\HS C=(\tr(C^*C))^{1/2}$, which makes each finite-dimensional matrix space a Hilbert space.

\begin{theorem}[Exterior Faber coefficients and the positive kernel]\label{lem:faber}
Let $\Omega$ be a bounded convex domain with regular real-analytic Jordan boundary. Let $h$ be meromorphic and one-to-one on $|\lambda|>r$ including infinity, where $0<r<1$, with a simple pole at infinity as in \eqref{eq:h}, and mapping $|\lambda|>1$ onto $\C\setminus\overline\Omega$. Put $q(\lambda)=\lambda h'(\lambda)$ and define $b_k(z)$ by expansion at infinity:
\begin{equation}\label{eq:faber}
 \frac{q(\lambda)}{h(\lambda)-z}=\sum_{k\ge0}b_k(z)\lambda^{-k}.
\end{equation}
Then $b_0=1$, and each $b_k$ is a polynomial of degree $k$ with leading coefficient $c^{-k}$. There are scalars $s_{kj}$, $k,j\ge1$, such that
\begin{equation}\label{eq:boundary-basis}
 b_k(h(\mu))=\mu^k+\sum_{j\ge1}s_{kj}\mu^{-j}\quad(k\ge1).
\end{equation}
The function
\begin{equation}\label{eq:s}
 s(\lambda,\mu)=\frac{q(\lambda)}{h(\lambda)-h(\mu)}-\frac{\lambda}{\lambda-\mu}
 =\sum_{k,j\ge1}s_{kj}\lambda^{-k}\mu^{-j}
\end{equation}
extends holomorphically across the diagonal and both infinities in $|\lambda|,|\mu|>r$. For every $r<\rho<1$ there is a constant $C_\rho$ with
$|s_{kj}|\le C_\rho\rho^{k+j}$. The series converges normally on
the exterior product collar: it converges absolutely and uniformly on
every compact subset, including subsets meeting either infinity.
In particular it converges absolutely and uniformly on $\T^2$. The continuous kernel
\[
 K(\lambda,\mu)=1+s(\lambda,\mu)+\overline{s(\lambda,\mu)}
\]
is nonnegative and has integral one in each variable with the other fixed.
\end{theorem}

\begin{proof}
Expanding the denominator in \eqref{eq:faber} after factoring $c\lambda$ shows that its $k$th coefficient is a polynomial of degree $k$ with leading term $c^{-k}z^k$. The constant coefficient is one. In particular, the $b_k$ form a basis of the scalar polynomials.

For the joint extension, consider
\[
 H(\lambda,\mu)=\frac{h(\lambda)-h(\mu)}{\lambda-\mu}.
\]
Write $u=\lambda^{-1}$, $v=\mu^{-1}$ and $h(\lambda)=c\lambda+g(u)$, where $g$ is holomorphic on $|u|<r^{-1}$. Then
\begin{equation}\label{eq:H}
 H=c-uv\frac{g(u)-g(v)}{u-v}.
\end{equation}
The divided difference of $g$ is jointly holomorphic, with diagonal value $g'(u)$, so this formula also covers either or both infinities. Off the diagonal and away from the axes, $H$ is nonzero by injectivity of $h$. On the diagonal it equals $h'(\lambda)\ne0$; on either axis it equals $c$. Thus it has no zeros on the reciprocal bidisk.

Direct differentiation gives $s=\lambda\partial_\lambda H/H=-u\partial_uH/H$. This function is holomorphic and vanishes on both reciprocal axes. Its Taylor series has precisely the strictly mixed powers in \eqref{eq:s}. For any $r<\rho<1$, Cauchy estimates on the closed reciprocal bidisk
of radius $\rho^{-1}$ give a constant $C_\rho$ with
$|s_{kj}|\le C_\rho\rho^{k+j}$. On a compact subset of the exterior
product collar, choose $\sigma>r$ below all finite coordinate moduli,
and then $r<\rho<\min(1,\sigma)$. The resulting geometric double
series dominates the Laurent series uniformly there; terms vanish on
the reciprocal axes. This proves normal convergence throughout the
product collar and absolute uniform convergence on the boundary. Comparing coefficients at $\lambda=\infty$ in \eqref{eq:s}, first for fixed finite $\mu$ in the exterior collar, yields \eqref{eq:boundary-basis}.

For distinct $\lambda,\mu\in\T$, one has $2\Re(\lambda/(\lambda-\mu))=1$, and therefore
\[
 K(\lambda,\mu)=2\Re\frac{q(\lambda)}{h(\lambda)-h(\mu)}\ge0
\]
by \eqref{eq:support}. Continuity gives the same sign on the diagonal. Every term of $s$ and of its conjugate has nonzero frequency in each variable. Their integrals in either variable vanish, proving both normalizations.
\end{proof}

To identify the basis, let $\Phi=h^{-1}$ in the exterior. Equation~\eqref{eq:boundary-basis}, which holds throughout the exterior collar, gives $b_k(z)-\Phi(z)^k=O(z^{-1})$ at infinity. Thus $b_k$ is the polynomial part of $\Phi^k$, the Faber polynomial associated with the chosen exterior coordinate. The usual normalization chooses the phase of $h$ so that $c>0$ \cite{Beckermann2005,BeckermannCrouzeix2014}.

These mixed coefficients are differentiated Grunsky coefficients.
The classical normalization gives a weighted coefficient inequality
\cite[equations~(1.1)--(1.3)]{BeckermannStylianopoulos2018};
the unweighted estimate used here follows from convexity and the two
angular marginal identities.

The two normalizations in Theorem~\ref{lem:faber} use the same exterior circle measure. Together they control the negative Fourier part of a boundary polynomial by its positive part, as follows.

\begin{proposition}[The coefficient norm estimate]\label{prop:coefficients}
For every positive integer $m$ and every $M_m(\C)$-valued polynomial $G(z)=\sum_{k=0}^N b_k(z)C_k$,
\begin{equation}\label{eq:coefficient-estimate}
 \sum_{k=0}^N\HS{C_k}^2
 \le\int_\T\HS{G(h(\mu))}^2\,d\tau(\mu)
 \le\HS{C_0}^2+2\sum_{k=1}^N\HS{C_k}^2.
\end{equation}
\end{proposition}
\begin{proof}
For $u\in L^2(\T;\mathcal H)$ with values in a Hilbert space $\mathcal H$, define
\[
 (Tu)(\mu)=\int_\T K(\lambda,\mu)u(\lambda)\,d\tau(\lambda).
\]
Positivity and the normalization in $\lambda$ imply
$\|(Tu)(\mu)\|^2\le\int K(\lambda,\mu)\|u(\lambda)\|^2d\tau(\lambda)$.
Integrating in $\mu$ and using the normalization in that variable shows that $T$ is an $L^2$ contraction.

For $u(\lambda)=\sum_{k=1}^N C_k\lambda^k$, only the $s$ term contributes to the integral, so
\begin{equation}\label{eq:kernel-fourier-action}
 (Tu)(\mu)=\sum_{j\ge1}\left(\sum_{k=1}^Ns_{kj}C_k\right)\mu^{-j}.
\end{equation}
All interchanges are justified by the absolute uniform convergence in Theorem~\ref{lem:faber}. In particular,
\[
 \sum_{j\ge1}\HS{\sum_{k=1}^Ns_{kj}C_k}^2\le\sum_{k=1}^N\HS{C_k}^2.
\]
By \eqref{eq:boundary-basis}, the boundary value of $G$ is the sum of its constant term $C_0$, its positive part $u$, and its negative part $Tu$. These have disjoint Fourier supports. Parseval's identity gives
\[
 \int_\T\HS{G\circ h}^2d\tau
 =\HS{C_0}^2+\sum_{k=1}^N\HS{C_k}^2+
   \sum_{j\ge1}\HS{\sum_{k=1}^Ns_{kj}C_k}^2,
\]
which proves both inequalities.
\end{proof}
For comparison with \cite[Proposition~3.1]{OpenAI2026Complete}, the
companion's exterior map $G$ and variables $t,w$ correspond to
$h,\lambda,\mu$, respectively.
The mixed correction is $s(\lambda,\mu)=B(\mu,\lambda)$, so
$s_{kj}=\beta_{jk}$ in that paper's notation. In particular,
\eqref{eq:kernel-fourier-action} uses this coefficient orientation and
does not require symmetry of the kernel.
\section{Singular vectors and the amplified estimate}\label{sec:matrix}
Fix $A\in M_n(\C)$ with $W(A)\subset\Omega$.
To estimate the norm of one polynomial evaluation, we form coefficient
functionals from its left and right singular vectors. The resolvent supplies
their coefficient arrays, while its positive Hermitian part compares the
resulting boundary norms. This Hermitian field is the numerical-range double-layer kernel of
\cite{DelyonDelyon1999,CrouzeixPalencia2017}; we prove the required
positivity for the original matrix $A$.

\begin{theorem}[The resolvent expansion and double-layer positivity]\label{lem:resolvent}
Let $\Omega,h,q$, and $(b_k)$ satisfy the hypotheses and conclusions of Theorem~\ref{lem:faber}. For any positive integer $n$ and $A\in M_n(\C)$ with $W(A)\subset\Omega$, on $\T$ let
\begin{equation}\label{eq:R}
 R(\lambda)=q(\lambda)(h(\lambda)I-A)^{-1}.
\end{equation}
Then
\[
 R(\lambda)=\sum_{k\ge0}b_k(A)\lambda^{-k},\qquad R(\lambda)+R(\lambda)^*\succeq0.
\]
For the fixed $A,h$, there are constants $C<\infty$ and $0<\theta<1$
such that $\|b_k(A)\|\le C\theta^k$ for all $k\ge0$. Thus the series
converges absolutely and uniformly on $\T$. In particular,
$\int_\T R\,d\tau=I$ and $\int_\T(R+R^*)\,d\tau=2I$.
\end{theorem}
\begin{proof}
Every eigenvalue of $A$ belongs to $W(A)$, by testing a unit eigenvector. Thus the exterior resolvent is defined on and outside the boundary. Compactness of the circle and strict spectral separation allow a smaller collar $|\lambda|>r_A$, with $r<r_A<1$, on which it remains holomorphic, including infinity where $R(\infty)=I$. Expansion at infinity gives the coefficients $b_k(A)$: this follows either by the geometric resolvent series for large $|h(\lambda)|$ or by substituting $A$ into the polynomial coefficient identities in \eqref{eq:faber}. Choose $r_A<\theta<1$. Cauchy estimates on the reciprocal circle of
radius $\theta^{-1}$ give the stated geometric bound. They imply
absolute uniform convergence on $\T$, and its constant coefficient
$b_0(A)=I$ gives both integrals.

For $B=h(\lambda)I-A$, direct congruence gives
\[
 B^*(R+R^*)B=\overline{q(\lambda)}B+q(\lambda)B^*.
\]
For each unit vector $v$, the quadratic form on the right is
$2\Re(\overline{q(\lambda)}(h(\lambda)-v^*Av))$, which is nonnegative by \eqref{eq:support}. Invertibility of $B$ proves positivity.
\end{proof}

\begin{theorem}[The bound on an analytic convex domain]\label{thm:domain}
For positive integers $n,m$, let $\Omega\subset\C$ be bounded, open, and convex with regular real-analytic Jordan boundary, and let $A\in M_n(\C)$ satisfy $W(A)\subset\Omega$. For every $M_m(\C)$-valued polynomial $F$ with $\|F(z)\|\le1$ on $\partial\Omega$, one has $\|F[A]\|\le2$.
\end{theorem}
\begin{proof}
Choose $h$ from Lemma~\ref{lem:exterior-collar} and apply
Theorem~\ref{lem:faber} to this coordinate.
Put $\gamma=\|F[A]\|$. There is nothing to prove if $\gamma=0$. Otherwise choose unit singular vectors $x,y\in\C^n\otimes\C^m$ such that
\begin{equation}\label{eq:singular}
 F[A]x=\gamma y,\qquad F[A]^*y=\gamma x.
\end{equation}
Write $x=\sum_{i=1}^mXe_i\otimes e_i$ and $y=\sum_{i=1}^mYe_i\otimes e_i$, where $X,Y$ are $n$-by-$m$ matrices and $(e_i)$ is the standard basis of $\C^m$. For the resolvent in Theorem~\ref{lem:resolvent}, define matrix functions
\[
 P=X^*RX,\quad Q=Y^*RY,\quad M=Y^*RX,\quad L=X^*RY.
\]
Subscripts $k\ge0$ denote the coefficients of $\lambda^{-k}$, so $M_k=Y^*b_k(A)X$, for example. By Theorem~\ref{lem:resolvent}, all four series converge absolutely and uniformly on $\T$, and their coefficient sequences are square summable. Since $b_0=1$,
\begin{equation}\label{eq:means}
 P_0=X^*X=P_0^*,\qquad Q_0=Y^*Y=Q_0^*,\qquad L_0^*=M_0.
\end{equation}
Set $a^2=\sum_{k\ge0}\HS{M_k}^2$. We will bound the squared
boundary $L^2$ Hilbert--Schmidt norms of $P+P^*$ and $Q+Q^*$
above by $4a^2/\gamma^2$ each, and that of $M+L^*$ below by $a^2$.
Positivity will then compare these quantities and force $\gamma\le2$.

\emph{The off-diagonal coefficient functional.}
For a polynomial $G=\sum_k b_kC_k$, direct expansion in the tensor basis gives
\begin{equation}\label{eq:pairing}
 y^*G[A]x=\sum_{k,i,j}(M_k)_{ij}(C_k)_{ij}.
\end{equation}
Indeed, $C_ke_j=\sum_i(C_k)_{ij}e_i$ in the coefficient factor, leaving $(Ye_i)^*b_k(A)Xe_j$ in the base factor. Equation~\eqref{eq:pairing} is the bilinear entry pairing of the two coefficient arrays. Cauchy--Schwarz gives
\begin{equation}\label{eq:crossfunctional}
 |y^*G[A]x|\le a\left(\sum_k\HS{C_k}^2\right)^{1/2}.
\end{equation}
Taking $G=F$ and using \eqref{eq:singular} shows that $a>0$.

\emph{The two ordered tests.}
Polynomial evaluation respects multiplication in its given order. Explicitly, if $F(z)=\sum_iD_iz^i$ and $G(z)=\sum_jE_jz^j$, then
\[
 (FG)[A]=\sum_{i,j}A^{i+j}\otimes D_iE_j=F[A]G[A],
\]
and the reversed product has coefficients $E_jD_i$. Hence
\begin{equation}\label{eq:ordered-products}
 y^*(FG)[A]x=\gamma x^*G[A]x,\qquad
 y^*(GF)[A]x=\gamma y^*G[A]y.
\end{equation}
Write the full finite Faber expansions as
\[
 FG=\sum_k b_kU_k,\qquad GF=\sum_k b_kV_k,
 \qquad \|H\|_\partial^2=\int_\T\HS{H(h(\mu))}^2\,d\tau(\mu).
\]
The functional in \eqref{eq:crossfunctional} uses all $M_k$, so it applies
to these full expansions for every finite degree. Apply it to the two
identities in \eqref{eq:ordered-products}. The lower coefficient estimate
in Proposition~\ref{prop:coefficients} and the two boundary multiplier
bounds give
\begin{align*}
 \gamma|x^*G[A]x|
 &\le a\Bigl(\sum_k\HS{U_k}^2\Bigr)^{1/2}
 \le a\|FG\|_\partial\le a\|G\|_\partial,\\
 \gamma|y^*G[A]y|
 &\le a\Bigl(\sum_k\HS{V_k}^2\Bigr)^{1/2}
 \le a\|GF\|_\partial\le a\|G\|_\partial.
\end{align*}
Here $\HS{F(z)G(z)}\le\|F(z)\|\HS{G(z)}$ and
$\HS{G(z)F(z)}\le\HS{G(z)}\|F(z)\|$ are used in their
respective orders. The upper coefficient estimate for $G$ now gives
\begin{equation}\label{eq:diagonal-functional}
 \max\{|x^*G[A]x|,|y^*G[A]y|\}
 \le\frac a\gamma
 \left(\HS{C_0}^2+2\sum_{k\ge1}\HS{C_k}^2\right)^{1/2}.
\end{equation}
The entries paired with $C_k$ in these two functionals are respectively those of $P_k$ and $Q_k$, by the calculation in \eqref{eq:pairing}. The weighted coefficient norm on the right gives reciprocal weights on these two coefficient sequences:
\begin{equation}\label{eq:weighted-dual}
 \HS{P_0}^2+\tfrac12\sum_{k\ge1}\HS{P_k}^2\le a^2/\gamma^2,
 \qquad
 \HS{Q_0}^2+\tfrac12\sum_{k\ge1}\HS{Q_k}^2\le a^2/\gamma^2.
\end{equation}
For the first inequality, set
$d_N=\HS{P_0}^2+\tfrac12\sum_{k=1}^N\HS{P_k}^2$
and test \eqref{eq:diagonal-functional} with $C_0=\overline{P_0}$,
$C_k=\overline{P_k}/2$ for $1\le k\le N$, and $C_k=0$ otherwise;
bars mean entrywise conjugation. The paired value is $d_N$, and the squared
weighted test norm is also $d_N$. Hence
$d_N\le(a/\gamma)\sqrt{d_N}$, which gives $d_N\le a^2/\gamma^2$
including when $d_N=0$. Letting $N\to\infty$ proves the first bound.
The same finite tests with $Q_k$ prove the second.

\emph{The boundary norms of the three blocks.}
The compression of $R+R^*$ by $[X\ Y]$ has diagonal blocks $P+P^*$,
$Q+Q^*$ and lower off-diagonal block $M+L^*$. We first compute their
boundary norms. In the diagonal fields, \eqref{eq:means} makes the
constant term twice $P_0$ or $Q_0$. For $k\ge1$, the coefficients $P_k$
and $P_k^*$ occur at frequencies $-k$ and $k$ and have equal
Hilbert--Schmidt norms; the same holds for $Q_k$ and $Q_k^*$.
Parseval therefore gives four times the reciprocal weighted sum in
\eqref{eq:weighted-dual}:
\begin{equation}\label{eq:diagonal-L2}
 \int_\T\HS{P+P^*}^2d\tau
 =4\left(\HS{P_0}^2+\tfrac12\sum_{k\ge1}\HS{P_k}^2\right)
 \le4a^2/\gamma^2,
\end{equation}
and the same identity and bound hold for $Q+Q^*$.

For the cross field, $L_0^*=M_0$ makes the constant term $2M_0$.
Its negative coefficients are $M_k$ for $k\ge1$ and its positive
coefficients are $L_k^*$ for $k\ge1$. Consequently
\begin{equation}\label{eq:cross-L2}
 \int_\T\HS{M+L^*}^2d\tau
 =4\HS{M_0}^2+\sum_{k\ge1}\bigl(\HS{M_k}^2+\HS{L_k}^2\bigr)
 \ge a^2.
\end{equation}

\emph{The positive block comparison.}
Theorem~\ref{lem:resolvent} gives $R+R^*\succeq0$. Its compression by the
$n$-by-$2m$ matrix $[X\ Y]$ is therefore
\[
 \Delta=\begin{pmatrix}P+P^*&L+M^*\\M+L^*&Q+Q^*\end{pmatrix}\succeq0.
\]
For $J=\diag(I_m,-I_m)$, the matrix $J\Delta J$ is also positive, so
\[
 0\le\tr\bigl(\Delta^{1/2}(J\Delta J)\Delta^{1/2}\bigr)
 =\tr(\Delta J\Delta J)
 =\HS{P+P^*}^2+\HS{Q+Q^*}^2-2\HS{M+L^*}^2.
\]
Integrating this inequality and using \eqref{eq:diagonal-L2} and \eqref{eq:cross-L2} yields
\[
 2a^2\le2\int_\T\HS{M+L^*}^2d\tau
 \le\int_\T\bigl(\HS{P+P^*}^2+\HS{Q+Q^*}^2\bigr)d\tau
 \le8a^2/\gamma^2.
\]
Since $a>0$, this proves $\gamma\le2$.
\end{proof}


% --- next excerpt ---

\begin{lemma}[Analytic convex outer approximation]\label{lem:approx}
Let $K\subset\C$ be nonempty, compact, and convex. There are bounded convex open domains $\Omega_j$ with regular real-analytic Jordan boundaries such that $K\subset\Omega_j$ and
\[
 \sup_{z\in\overline\Omega_j}\dist(z,K)\longrightarrow0.
\]
\end{lemma}
